\documentclass[11pt]{article}
\usepackage[letterpaper,margin=1in]{geometry}
\usepackage{graphicx}
\usepackage{hyperref}

\usepackage[backend=biber,
  maxbibnames=99,
  natbib=true,
  style=ext-numeric-comp,
%  autocite=numeric, %superscript,
  sortcites=true,
  sorting=nyt]{biblatex}
\DeclareFieldFormat % Titles in sentence case
  [article,inbook,incollection,inproceedings,patent,thesis,unpublished]
  {titlecase:title}{\MakeSentenceCase*{#1}}
\renewbibmacro{in:}{ % Don't use "In:" before journal names.  
  \ifentrytype{article}{}{\printtext{\bibstring{in}\intitlepunct}}}
% date after authors
\renewbibmacro*{date+extradate}{}
\renewbibmacro*{issue+date}{}
\renewbibmacro*{date+extradate}{%
  \printtext[parens]{%
    \printdate
  }%
}
\renewbibmacro*{author}{%
  \ifuseauthor
    {\printnames{author}%
     \setunit{\addspace}%
     \usebibmacro{date+extradate}}
    {}%
}
\renewcommand*{\multicitedelim}{\addsemicolon\space}

\addbibresource{defs.bib}
\addbibresource{arrpubs.bib}
\addbibresource{arr.bib}
\addbibresource{molrec.bib}
\addbibresource{qchar.bib}

\renewcommand{\topfraction}{.9}
\renewcommand{\bottomfraction}{.9}
\renewcommand{\textfraction}{.1}
\renewcommand{\floatpagefraction}{.8}
\renewcommand{\dbltopfraction}{.9}
\renewcommand{\dblfloatpagefraction}{.9}
\setcounter{topnumber}{2}
\setcounter{bottomnumber}{2}
\setcounter{totalnumber}{5}
\setcounter{dbltopnumber}{2}
\setcounter{secnumdepth}{0}
\begin{document}
\title{Notes on Genetic Group Differences}
\author{Alan R. Rogers}
\maketitle

Human racial classification is of no social value and is positively
destructive of social and human relations. Since such racial
classifcation [sic] is now seen to be of virtually no genetic or
taxonomic significance either, no justification can be offered for its
continuance. (Lewontin, 1972, final paragraph)

Yair and Coop [18] (2023) consider the case in which stabilizing
selection operates with identical fitness functions in several
populations. Stabilizing selection prevents the means of the
quantitative character from diverging, so the populations retain
similar character means. But at each contributing locus, selection
favors whichever allele is common. Consequently, an allele that drifts
past the minimum of this marginal (1-locus) fitness function is then
favored by selection. This can accelerate the divergence of polygenic
scores for the quantitative character. Variation among these scores
can exceed neutral expectations.

Look at Shen and Feldman [6]

``Lewontin recognized the importance of genetic variation among
populations in biomedical problems; indeed, Feldman & Lewontin [43]
included a perspective on the value of genetic ancestry in such
contexts.''

``Lewontin’s brief concluding interpretation is incomplete. A more
comprehensive view can be found by studying the 1972 paper together
with Lewontin’s other writings, as in Shen & Feldman [6].'' p. 5

5: ``The similarity of biological-racist claims in the time of
Lewontin’s paper and in the present are apparent in the references
tabulated by Carlson \& Harris [7]''

Shen and Feldman 3: ``If Mitton and Edwards were merely\ldots arguing
that human genetic variation can be used to establish a useful
taxonomy, we believe that Lewontin would agree, as expressed in
Feldman & Lewontin [34]. However, Mitton and Edwards supported the
potential utility of the classical race classification, and therein
lies the true problem with their critique of Lewontin.''

Shen and Feldman 3: ``It is true that small differences in allele
frequencies at large numbers of markers, with appropriate algorithms
[39–41], allow continental clustering [15,20,21]. For pairs of
populations from different clusters, genetic distance is generally
larger than between pairs of populations from the same clusters,
indicating that the history of migration is important
[34,42,43]. However, for Lewontin, ‘racial classification’ had a
totally different meaning from continental origin.''

Me: Shen and Feldman presume to tell us what Lewontin meant by the
word ``race''. According to them, Lewontin's conclusion meant only that
one particular classification was wrong and left open the possibility
that other classifications might work better. If that were the case,
then Lewontin would presumably have embraced the more recent
classifications produced from molecular data. But Lewontin never did
that, although he lived long enough to. Neither have Shen or
Feldman. This suggests that Lewontin, Shen, and Feldman were not
making a narrow claim about one particular classification of humans
into races. They were making the broad claim that no classification of
humans is possible. This is manifestly untrue, and Edwards was
therefore correct in calling Lewontin's conclusion false.

Shen and Feldman 5: ``Analyses of genetic and phenotypic
differentiation do not exclude the possibility that some traits
(possibly viewed as ‘essential’) exhibit a similar pattern [of large
between-group differences] to that of skin colour.''

Shen and Feldman 5: Lewontin [4, 6] is a critique of heritability

Edwards 799: ``Now suppose there are k similar loci, all with gene
frequency p in population 1 and q in population 2. The ratio of the
within-to-total variability is still 84% at each locus. The total
number of ‘þ’ genes in an individual will be binomial with mean kp in
population 1 and kq in population 2, with variance kpq in both
cases. Continuing with the former gene frequencies and taking k¼ 100
loci (say), the mean numbers are 30 and 70 respectively, with
variances 21 and thus standard deviations of 4.58. With a difference
between the means of 40 and a common standard deviation of less than
4.6, there is virtually no overlap between the distributions.''

Appiah:CI-12-21: ``The criterion Du Bois actually uses amounts to
this: people are members of the same race if they share features in
virtue of being descended largely from people of the same region.''

Appiah:CI-12-21, p.~37, footnote 10, commenting on Nei and
Roychoudhury Genetic Relationship and Evolution of Human Races,
Evolutionary Biology 14 (1983): 11.: ``I would dispute their claim
that their work shows the existence of a biological basis for the
classification of human races; what it shows is that human populations
differ in their distributions of genes. That is a biological fact. The
objection to using this fact as a basis of a system of classification
is that far too many people don't fit into just one category that can
be so defined.''

Roseman:BE-43-2100204, p3: Roseman's concern seems to be not that
evolutionary psychologists are comparing one race to another, but
rather that their race categories are incorrect--that they differ
from those implied by genetic data. ``Suppose we apply our racial
model (Figure 1C,D) to our final scenario (Figure 1E,F) as some
evolutionary psychologists advocate. We would be ignoring key
components of the correlations in evolutionary outcomes, such as the
substantial correlation between group B and groups C and
D. Likewise, we would be inserting a positive correlation between A
and B where none indeed exists. The mistake is further compounded by
the fact that gene flow and different rates of evolution by drift
are important causes of the correlation structure of human
variation.[3,4,7] This is intentional ignorance of the correlation
structure of the data, which is Lewontin’s fallacy.''

Roseman:BE-43-2100204, p4: ``Contrary to Edwards’ claims, Lewontin
did consider the correlation structure that was apparent in the data
as his estimates of among race variation depend on nonzero
correlations among evolutionary outcomes at the racial level.'' He
says this because Lewontin's data (had they been an analysis of
variance instead of entropy) can be used to calculate the intraclass
correlation of individuals within subdivisions. But Edwards was
talking about the correlation among loci (between-group LD), not the
correlation among individuals.

Shen and Feldman assert that when Lewontin spoke of ``races'', he did
not mean ``continental human populations''. But here is what Lewontin
himself said: ``only 7 percent of all human genetic variation lies on
the average between major human races such as those of Africa, Asia,
Europe, and Oceania.'' This quote is from Pearce et al (doi:
10.1136/bmj.328.7447.1070) who were quoting from Lewontin (1991)
``Biology as Ideology: The Doctrine of DNA.''

VanArsdale:ARA-48-227, p 228: ``variation within populations exceeds
the variation that serves to distinguish between groups of
populations, discrediting the significance of the latter.''

VanArsdale:ARA-48-227, p 230: ``Small differences do not presuppose
insignificant differences.''

Koenig:RRG-08-1, p 1: ``for decades, the finding that there is
greater genetic variability `within groups' than `between groups'
had generally been accepted as evidence that the human species is
not divided into discrete races.''

Risch:GB-3-2007.1, p.~2007.4: ``Racial categorizations have never
been based on skin pigment, but on indigenous continent of origin.''

Sankar:RRG-08-271, p.~275: Cites Gannet (2003) and Sauer (1993) re
the views of ``early typological race theorists''.

Long:HB-81-777, p.~795: ``The relative proportion of variation within
and among groups therefore appears to be meaningless as a criterion
for judging the validity of races or subspecies as defined by
biologists.''

Long:HB-81-777, p.~796: ``Four decades ago, Frank Livingstone
declared the nonexistence of human races (Livingstone 1963). It is
now time for geneticists and anthropologists to stop worrying about
what does not exist and to discover what does exist.''

Crow:D-131-88, p.~84: ``Unlike those anthropologists who deny the
usefulness of the term, I believe that the word `race' can be
meaningfully applied to groups that are partially mixed.''

Crow:D-131-88, p.~85: ``Far more contentious is the evidence that
some skills and behavioral properties are differentially distributed
among different racial groups. There is strong evidenence that such
racial differences are partly genetic, but the evidence is more
indirect and has not been convincing to everyone.''

Lewontin:NOG-84, p.~126: ``Of all human genetic variation known for
enzymes and other proteins, where it has been possible to actually
count up the frequencies of different forms of the genes and so to
get an estimate of genetic variation, 85 percent turns out to be
between individuals within the same local population, tribe, or
nation; a further 8 percent is between tribes or nations within a
major `race'; and the remaining 7 percent is between major `races.'\,''

Lewontin:NOG-84, p.~127: ``The differences between major `racial'
categories, no matter how defined, turn out to be rather
small. Human `racial' differentiation is, indeed, only skin
deep. Any use of racial categories must take its justifications from
some other source than biology.''

Turner:PLO-3-e285, p.~1572: Genomic islands of speciation ``We found
that three regions, together encompassing less than 2.8 Mb, are the
only locations where the M and S forms are significantly
differentiated. Two of these regions are adjacent to centromeres, on
Chromosomes 2L and X, and contain 50 and 12 predicted genes,
respectively. Sequenced loci in these regions contain fixed
differences between forms and no shared polymorphisms, while no
fixed differences were found at nearby control loci. The third
region, on Chromosome 2R, contains only five predicted genes; fixed
differences in this region were also verified by direct
sequencing. These `speciation islands' remain differentiated despite
considerable gene flow, and are therefore expected to contain the
genes responsible for reproductive isolation.''

Turner:PLO-3-e285, p.~1572: ``A. gambiae is further subdivided into
two partially isolated taxa known as the M form and S form
[18]. \ldots Subsequent studies of reproductive isolation in nature
revealed that these forms mate assortatively, with 98.8\% of
wild-caught gravid females (within an area of sympatry) having mated
within their own form [20,21]. When forms are crossed in the lab, no
intrinsic fitness reductions are found [22], suggesting that the
observed heterozygotic deficiencies are due to nonrandom mating
and/or ecologically dependent postzygotic isolation.'' 

Turner:PLO-3-e285, p.~1576: ``A previous survey using microsatellite
loci of M and S forms in Cameroon found Fst values between forms
that were consistent with substantial migration rates (calculating
4Nem from the average Fst reported in [26] yields $10 < 4N_e m <
47$).'' Using $F_{ST} = 1/(1 + 4N_e m)$, this corresponds to $0.02 <
F_{ST} < 0.09$.

Turner:PLO-3-e285, p.~1576: ``In contrast to the low levels of
differentiation found throughout most of the genome, our array
experiments revealed three small regions to be significantly
differentiated between forms. Sequences from islands on Chromosomes
2L and X contain 13 fixed differences and no shared
polymorphisms. One gene within the third island, a 37-kb region on
Chromosome 2R, also shows only fixed differences between forms.''

Turner:PLO-3-e285, p.~1576: ``We conducted coalescent simulations to
test whether the observed differentiation on Chromosomes 2L and X
could plausibly result from neutral scenarios. Rejection of this
neutral hypothesis ($p < 0.0001$) suggests that differentiation in
these regions is due to selection against hybrid genotypes during
backcrossing. This conclusion supports the prediction that when gene
flow is present, differentiation between incipient species can be
limited to small regions surrounding isolating genes [37].''

Turner:ME-19-848, p~849: ``With the exception of rare hybrid
individuals (more about this below), sequenced loci in these
pericentromeric regions show fixed DNA differences between forms,
with no polymorphisms shared between them. This is in stark contrast
to the remaining 97\% of the assembled genome, where fixed
differences have not yet been found, and many polymorphisms are
shared. Though allele frequency differences in the rest of the
genome are often adequate to cluster samples based on their M ⁄ S
designation, this degree of differentiation is slight compared to
fixed differentiation seen in the pericentromeric regions.''

Turner:ME-19-848, p~849: ``A pattern of low $F_{ST}$ values across
the majority of the genome, coupled with direct evidence of
hybridization in nature, would seem to support models of `mosaic’
speciation, where ongoing hybridization homogenizes variation in
most of the genome while loci under strong selection remain in
disequilibrium with each other (Wu & Ting 2004; Via & West 2008;
Nosil et al. 2009).''

Turner:ME-19-848, p~849: ``A reasonable alternative hypothesis—raised
by White and colleagues—is that there is actually little gene flow
between the M- and S-forms, and that the low level of
differentiation observed on chromosome arms is due to segregating
ancestral variation. Under this scenario, the hybrids observed in
nature would be mostly F1 individuals who die before reproducing. If
F1 genotypes are effectively sterile due to ecological or
behavioural maladaptation, then the amount of `realized’ gene flow
could be near zero. The divergent islands could appear different
from the rest of the genome because linked selection in these low
recombination regions reduces the effective population size, leading
to faster sorting of ancestral variation by fixing alternative
haplotypes in each form.''

Nosil:ME-18-375 p.~389: ``Five genome scan studies (1, 2, 4, 8, 9 in
Table 2) have analogously compared population trees for data sets
that either include or exclude outliers. All these studies observed
that outlier loci group populations more as a function of ecology,
whereas putatively neutral (nonoutlier) loci group populations in a
manner more consistent with geography (see also Via & West 2008).''

Nosil:ME-18-375 p.~389: ``The second example treats Helianthus annuus
and H. petiolaris sunflowers (Yatabe et al. 2007). These species
hybridize extensively, F1 hybrids exhibit extremely low pollen
fertility, and these fertility barriers map to chromosomal
rearrangements (Rieseberg et al. 1999). This study, however,
detected no association of outliers with QTL for morphological
differences or hybrid sterility, and only weak associations of
outliers with chromosomal inversions. The authors conclude that
regions of differentiation between these species are very small.''

Nosil:ME-18-375 p.~393: ``The chance that a new mutation (whether
adaptive or neutral) will persist and increase in frequency is
highest in regions of reduced introgression (Gavrilets 2004, chapter
6). Therefore, differentiated loci are expected to accumulate in
genomic regions that already harbour genes under divergent
selection, leading to increases in the number of genes within an
island, and thus, in island size, and potentially generating a
positive feedback loop. This model thus predicts that genes
affecting local adaptation will form clusters within the genome
rather than being more evenly distributed across it.''

Nosil:ME-18-375 p.~395: ``Divergent selection is predicted to yield
heterogeneous divergence across the genome.

Cruickshank:ME-23-3133, p.~3148: The main feature of the alternative
model---variously referred to as the `incidental island’ model (Turner
\& Hahn 2010), the `selection at linked sites’ model (Nachman \&
Payseur 2012) or the `low-gene-flow’ model (Hahn et al. 2012)---is
that heterogeneous levels of relative divergence are largely due to
heterogeneous natural selection across the genome and not
heterogeneous migration rates (Fig. 5b). Because any type of selection
that lowers levels of linked neutral diversity will result in
increased relative divergence (Charlesworth et al. 1997; Charlesworth
1998; Slatkin \& Wiehe 1998), differences in the effects of selection
across the genome will result in different levels of relative
divergence

Two populations with allele frequencies $p+\delta$ and
$p-\delta$. Probability of identity for two gene copies, one from each
population, is
\[
J = (p+\delta)(p-\delta) + (q-\delta)(q+\delta)
= p^2-\delta^2 + q^2 - \delta^2 = p^2 + q^2 - 2\delta^2
\]
Probability that they differ is
\[
d = 1-J = 2pq + 2\delta^2 = 2pq(1+F_{ST})
\]
where $F_{ST} = \delta^2/pq$. Divergent selection would tend to
increase $2pq$, because it would push allele frequency toward 1 in one
pop but toward 0 in the other, pushing the mean toward 1/2. Thus, $d$
should increase faster than $F_{ST}$.


\printbibliography

\end{document}
